\documentclass[aspectratio=169,10pt]{beamer}
% Shared Beamer setup for Fall 2026 course slides.
% Clean Cal Poly-inspired theme: no custom class files or uncommon packages.
\mode<presentation>{
  \usetheme{default}
  \usefonttheme{professionalfonts}
}

\definecolor{CalPolyGreen}{HTML}{154734}
\definecolor{CalPolyGold}{HTML}{BD8B13}
\definecolor{CalPolySage}{HTML}{7FA28F}
\definecolor{CalPolyMint}{HTML}{DDF1E7}
\definecolor{CalPolySand}{HTML}{A29F86}
\definecolor{CalPolyBeige}{HTML}{EFEEDE}
\definecolor{DeckInk}{HTML}{1F2933}
\definecolor{DeckMuted}{HTML}{5F6B7A}
\definecolor{DeckSoft}{HTML}{F7F7F5}

\setbeamersize{text margin left=0.55in,text margin right=0.55in}
\setbeamercolor{normal text}{fg=DeckInk,bg=white}
\setbeamercolor{structure}{fg=CalPolyGreen}
\setbeamercolor{title}{fg=white}
\setbeamercolor{frametitle}{fg=CalPolyGreen,bg=white}
\setbeamercolor{block title}{bg=CalPolySage,fg=white}
\setbeamercolor{block body}{bg=CalPolyMint,fg=black}
\setbeamercolor{alerted text}{fg=CalPolyGold}

\setbeamerfont{title}{size=\Huge,series=\bfseries}
\setbeamerfont{frametitle}{size=\Large,series=\bfseries}
\setbeamerfont{block title}{size=\normalsize,series=\bfseries}

\setbeamertemplate{navigation symbols}{}
\setbeamertemplate{itemize item}{\color{CalPolyGold}$\blacktriangleright$}
\setbeamertemplate{itemize subitem}{\color{CalPolyGreen}$\bullet$}
\setbeamertemplate{footline}{
  \hfill{\color{DeckMuted}\scriptsize\insertframenumber/\inserttotalframenumber}
  \hspace{0.18in}\vspace{0.12in}
}
\setbeamertemplate{frametitle}{
  \vspace{0.18in}
  \hspace{0.02in}\insertframetitle\par
  \vspace{0.08in}
  \noindent\makebox[\linewidth][l]{%
    {\color{CalPolyGold}\rule{0.55in}{2pt}}%
    {\color{CalPolyGreen}\rule{\dimexpr\linewidth-0.55in\relax}{2pt}}%
  }
}

\newcommand{\CourseNumber}{Course}
\newcommand{\CourseTitle}{Course Title}
\newcommand{\LectureNumber}{0}
\newcommand{\LectureTitle}{Lecture Title}
\newcommand{\LectureDate}{Date}
\newcommand{\CourseUnit}{Unit}

\newcommand{\coursenumber}[1]{\renewcommand{\CourseNumber}{#1}}
\newcommand{\coursetitle}[1]{\renewcommand{\CourseTitle}{#1}}
\newcommand{\lecturenumber}[1]{\renewcommand{\LectureNumber}{#1}}
\newcommand{\lecturetitle}[1]{\renewcommand{\LectureTitle}{#1}}
\newcommand{\lecturedate}[1]{\renewcommand{\LectureDate}{#1}}
\newcommand{\courseunit}[1]{\renewcommand{\CourseUnit}{#1}}

\author{Kirk Duran}
\institute{California Polytechnic State University, San Luis Obispo}
\date{\LectureDate}

\newcommand{\makedecktitle}{
  \begingroup
  \setbeamercolor{background canvas}{bg=CalPolyGreen}
  \begin{frame}[plain]
    \vfill
    \begin{center}
      {\usebeamerfont{title}\color{white}\LectureTitle\par}
      \vspace{0.18in}
      {\Large\color{white}Lecture \LectureNumber\par}
    \end{center}
    \vfill
    \noindent{\color{white}\rule{\linewidth}{1.2pt}}\par
    \vspace{0.08in}
    \noindent\colorbox{CalPolyGold}{\makebox[0.24\linewidth][c]{\color{white}\Large\CourseNumber}}
    \hspace{0.12in}{\color{white}\Large Kirk Duran}
  \end{frame}
  \endgroup
}

\newenvironment{keyidea}{\begin{block}{Key idea}}{\end{block}}
\newenvironment{activity}{\begin{block}{In class}}{\end{block}}
\newenvironment{cpdefinition}{\begin{block}{Definition}}{\end{block}}
\newenvironment{exampleblockcp}{
  \begingroup
  \setbeamercolor{block title}{bg=CalPolySand,fg=white}
  \setbeamercolor{block body}{bg=CalPolyBeige,fg=black}
  \begin{block}{Example}
}{
  \end{block}
  \endgroup
}

\newcommand{\imageplaceholder}[2][0.3in]{%
  \noindent\makebox[\linewidth][c]{%
    \fcolorbox{CalPolySage}{CalPolyBeige}{%
      \parbox[c][#1][c]{0.86\linewidth}{%
        \centering
        {\color{DeckMuted}\scriptsize Image placeholder: }%
        {\color{DeckInk}\scriptsize\ttfamily\detokenize{#2}}%
      }%
    }%
  }%
}

\newcommand{\sectionframe}[1]{
  \begingroup
  \setbeamercolor{background canvas}{bg=CalPolyGreen}
  \begin{frame}[plain]
    \vfill
    {\color{white}\LARGE\bfseries #1\par}
    \vspace{0.18in}
    {\color{CalPolyGold}\rule{0.22\paperwidth}{3pt}}
    {\color{white}\rule{0.62\paperwidth}{3pt}}
    \vfill
  \end{frame}
  \endgroup
}

\newcommand{\placeholdernote}{
  \begin{block}{Placeholder}
    Replace these bullets with the final lecture material, examples, and in-class activities.
  \end{block}
}

\usepackage[T1]{fontenc}
\usepackage{lmodern}
\usepackage{graphicx}
\usepackage{fancyvrb}
\usepackage{tabularx}
\graphicspath{{../assets/lecture-35/}}
\setbeamersize{text margin left=0.42in,text margin right=0.42in}
\newcommand{\lectureimage}[2][1.81in]{\includegraphics[width=0.93\linewidth,height=#1,keepaspectratio]{#2}}
\coursenumber{CS 1001}
\coursetitle{Introduction to Computing}
\lecturenumber{35}
\lecturetitle{Exceptions and Error Handling}
\lecturedate{November 20, 2026}
\courseunit{13. Input, Files, and Robust Programs}
\begin{document}
\makedecktitle
\begin{frame}[fragile]{Today: errors are part of the interface}
\begin{itemize}\item Recognize and handle common exceptions with \texttt{try} / \texttt{except}.\item Choose specific handlers, inspect error messages, and avoid over-catching.\item Understand propagation, \texttt{raise}, and custom exceptions.\item Verify failures with \textbf{two forms} of \texttt{unittest.assertRaises}.\end{itemize}
\vspace{0.06in}\begin{keyidea}\small Design for bad input, not just the happy path.\end{keyidea}
\end{frame}
\begin{frame}[fragile]{A file can exist and still contain bad data}
\begin{columns}[T,onlytextwidth]\begin{column}{0.52\textwidth}\begin{block}{Example}
\begin{Verbatim}[fontsize=\footnotesize]
line = "Ada,not-a-number"
name, raw_score = line.split(",")
score = int(raw_score)  # ValueError
\end{Verbatim}
\end{block}\end{column}\begin{column}{0.45\textwidth}\small\begin{itemize}\item Lecture 34 taught us how to read and parse text.\item But \texttt{int()} cannot convert arbitrary text.\item Should the whole program crash, ignore it, or explain the problem?\end{itemize}\end{column}\end{columns}
\vspace{0.06in}\begin{keyidea}\small Today: decide where and how failures are handled.\end{keyidea}
\end{frame}
\begin{frame}[fragile]{Two paths through one program}
\begin{center}\lectureimage{exception-fork.png}\end{center}
\small\begin{itemize}\item Normal flow produces a value.\item Exceptional flow produces an exception object.\item A matching handler can recover or report the problem.\end{itemize}
\end{frame}
\begin{frame}[fragile]{Definition: an exception}
\begin{itemize}\item An \textbf{exception} is an object signaling that an operation could not proceed normally.\item Python \textbf{raises} an exception when it detects a problem.\item An \texttt{except} clause can \textbf{catch} a matching exception.\item Without a handler, Python reports an unhandled exception and traceback.\end{itemize}
\vspace{0.06in}\begin{keyidea}\small An exception is not the same thing as printing an error message.\end{keyidea}
\end{frame}
\begin{frame}[fragile]{Python can raise exceptions for us}
\begin{block}{Python / sample output}
\begin{Verbatim}[fontsize=\footnotesize]
int("not a number")      # ValueError
10 / 0                   # ZeroDivisionError
values = [1, 2]
print(values[9])         # IndexError
open("missing.txt", "r") # FileNotFoundError
\end{Verbatim}
\end{block}
\small\begin{itemize}\item The \textbf{operation} decides when it cannot continue.\item Different failures have different exception types.\end{itemize}
\end{frame}
\sectionframe{Part 1: Catching Exceptions}
\begin{frame}[fragile]{Catch the conversion error}
\begin{block}{Python / sample output}
\begin{Verbatim}[fontsize=\footnotesize]
text = input("Enter a whole number: ")

try:
    number = int(text)
    print(number * 2)
except ValueError:
    print("Please enter a valid integer.")
\end{Verbatim}
\end{block}
\small\begin{itemize}\item Only a matching \texttt{ValueError} runs this handler.\item Execution resumes after the \texttt{try} / \texttt{except} statement.\end{itemize}
\end{frame}
\begin{frame}[fragile]{The \texttt{try} / \texttt{except} decision}
\begin{center}\lectureimage{try-except-control.png}\end{center}
\small\begin{itemize}\item No exception: skip the handler.\item Exception: jump to the first matching handler.\item Statements after the failing operation are skipped.\end{itemize}
\end{frame}
\begin{frame}[fragile]{A failing statement skips the rest of \texttt{try}}
\begin{block}{Python / sample output}
\begin{Verbatim}[fontsize=\footnotesize]
try:
    score = int("oops")
    print("converted")
except ValueError:
    print("invalid score")

print("afterwards")
\end{Verbatim}
\end{block}
\small\begin{itemize}\item Predict the printed lines.\item Why does \texttt{"converted"} never appear?\end{itemize}
\begin{activity}\small Ask students to trace the control flow before revealing the answer.\end{activity}
\end{frame}
\begin{frame}[fragile]{Output and reasoning}
\begin{block}{Python / sample output}
\begin{Verbatim}[fontsize=\footnotesize]
invalid score
afterwards
\end{Verbatim}
\end{block}
\small\begin{itemize}\item \texttt{int("oops")} raises \texttt{ValueError}.\item Control jumps straight to \texttt{except ValueError}.\item The program continues after the handled exception.\end{itemize}
\end{frame}
\begin{frame}[fragile]{What if the exception type does not match?}
\begin{block}{Python / sample output}
\begin{Verbatim}[fontsize=\footnotesize]
try:
    result = 10 / 0
except ValueError:
    print("Bad numeric text")
\end{Verbatim}
\end{block}
\small\begin{itemize}\item The error is \textbf{ZeroDivisionError}, not \texttt{ValueError}.\item No matching handler exists here.\item The exception propagates to the caller.\end{itemize}
\vspace{0.06in}\begin{keyidea}\small Catching one exception does not automatically catch every exception.\end{keyidea}
\end{frame}
\begin{frame}[fragile]{Handle several errors differently}
\begin{block}{Python / sample output}
\begin{Verbatim}[fontsize=\footnotesize]
try:
    number = int(raw)
    result = 100 / number
except ValueError:
    print("Expected an integer")
except ZeroDivisionError:
    print("Cannot divide by zero")
\end{Verbatim}
\end{block}
\small\begin{itemize}\item One operation may fail in different ways.\item Each branch gives a response appropriate to the failure.\end{itemize}
\end{frame}
\begin{frame}[fragile]{Handler order matters}
\begin{center}\lectureimage{specific-first-handlers.png}\end{center}
\small\begin{itemize}\item Python checks handlers from top to bottom.\item The \textbf{first} matching handler runs.\item Put specific exception types before a broad parent type.\end{itemize}
\end{frame}
\begin{frame}[fragile]{Combine types when the response is identical}
\begin{block}{Python / sample output}
\begin{Verbatim}[fontsize=\footnotesize]
try:
    answer = 100 / int(raw)
except (ValueError, ZeroDivisionError):
    print("Cannot compute an answer")
\end{Verbatim}
\end{block}
\small\begin{itemize}\item One handler can catch a \textbf{tuple} of exception types.\item Use separate branches if users need distinct messages.\end{itemize}
\end{frame}
\begin{frame}[fragile]{A generic catch: \texttt{Exception}}
\begin{block}{Python / sample output}
\begin{Verbatim}[fontsize=\footnotesize]
try:
    process_record(raw)
except Exception as exc:
    print("Something went wrong:", exc)
\end{Verbatim}
\end{block}
\small\begin{itemize}\item \texttt{Exception} catches most ordinary application errors.\item This is sometimes useful at the program boundary.\item Inside a helper, it can hide programming bugs.\end{itemize}
\vspace{0.06in}\begin{keyidea}\small Prefer the narrowest exception that you actually know how to handle.\end{keyidea}
\end{frame}
\begin{frame}[fragile]{Why not just write bare \texttt{except:}?}
\begin{itemize}\item \texttt{except:} also catches exceptions such as \texttt{KeyboardInterrupt} and \texttt{SystemExit}.\item It may stop a user from interrupting a program normally.\item \texttt{except Exception:} is narrower, but still broad.\item Often the correct choice is a specific type such as \texttt{ValueError} or \texttt{OSError}.\end{itemize}
\vspace{0.06in}\begin{keyidea}\small Do not silently swallow all failures with \texttt{except: pass}.\end{keyidea}
\end{frame}
\begin{frame}[fragile]{Capture the exception message}
\begin{block}{Python / sample output}
\begin{Verbatim}[fontsize=\footnotesize]
try:
    value = int("four")
except ValueError as exc:
    print("Conversion failed:", exc)
    print(type(exc).__name__)
\end{Verbatim}
\end{block}
\small\begin{itemize}\item \texttt{as exc} binds the \textbf{exception object}.\item \texttt{str(exc)} gives a readable message.\item The type and message can help explain or log the failure.\end{itemize}
\end{frame}
\begin{frame}[fragile]{Keep the \texttt{try} block small}
\begin{block}{Python / sample output}
\begin{Verbatim}[fontsize=\footnotesize]
try:
    count = int(raw_count)
except ValueError:
    print("Invalid count")
else:
    print("Count is", count)
\end{Verbatim}
\end{block}
\small\begin{itemize}\item Catch only failures you expect from the relevant operation.\item \texttt{else} executes only if the \texttt{try} block completed without an exception.\end{itemize}
\end{frame}
\begin{frame}[fragile]{Cleanup: \texttt{finally}, and why \texttt{with} still matters}
\begin{block}{Python / sample output}
\begin{Verbatim}[fontsize=\scriptsize]
try:
    print("start")
    value = int(raw)
finally:
    print("cleanup step")

# For files, prefer:
with open("scores.txt") as f:
    text = f.read()
\end{Verbatim}
\end{block}
\small\begin{itemize}\item \texttt{finally} runs when leaving \texttt{try} (including exceptional exits).\item Use \texttt{with} for files: it manages closing them.\item Do not replace \texttt{with open(...)} with manual close logic.\end{itemize}
\end{frame}
\sectionframe{Part 2: Propagation and \texttt{raise}}
\begin{frame}[fragile]{An exception travels up the call stack}
\begin{center}\lectureimage{stack-unwinding.png}\end{center}
\small\begin{itemize}\item \texttt{parse\_score()} can raise.\item \texttt{load\_scores()} can leave that exception unhandled.\item \texttt{main()} can catch it and decide what to tell the user.\end{itemize}
\end{frame}
\begin{frame}[fragile]{Propagation needs no special syntax}
\begin{block}{Python / sample output}
\begin{Verbatim}[fontsize=\footnotesize]
def parse_number(raw: str) -> int:
    return int(raw)

def double(raw: str) -> int:
    number = parse_number(raw)
    return number * 2

double("oops")  # ValueError propagates
\end{Verbatim}
\end{block}
\small\begin{itemize}\item No \texttt{except} in either helper.\item The caller receives the raised exception automatically.\end{itemize}
\end{frame}
\begin{frame}[fragile]{Handle at the right boundary}
\begin{block}{Python / sample output}
\begin{Verbatim}[fontsize=\scriptsize]
def parse_number(raw: str) -> int:
    return int(raw)

def main() -> None:
    raw = input("Number: ")
    try:
        print(parse_number(raw))
    except ValueError:
        print("Please enter a valid number")
\end{Verbatim}
\end{block}
\small\begin{itemize}\item Helpers report failure through exceptions.\item The user-facing \texttt{main()} chooses the friendly message.\end{itemize}
\end{frame}
\begin{frame}[fragile]{Raise a problem deliberately}
\begin{block}{Python / sample output}
\begin{Verbatim}[fontsize=\footnotesize]
def set_age(age: int) -> None:
    if age < 0:
        raise ValueError("age cannot be negative")
    print("Age accepted:", age)
\end{Verbatim}
\end{block}
\small\begin{itemize}\item Use \texttt{raise ExceptionType("message")}.\item Code after \texttt{raise} does not run on that path.\item The caller may catch the exception or let it propagate.\end{itemize}
\end{frame}
\begin{frame}[fragile]{When is \texttt{None} better than an exception?}
\begin{center}\lectureimage{return-vs-raise.png}\end{center}
\small\begin{itemize}\item \texttt{None} can mean an ordinary, expected absence.\item An exception can represent an invalid input that violates a function contract.\item Be consistent: document what callers should expect.\end{itemize}
\end{frame}
\begin{frame}[fragile]{The difference between \texttt{raise} and \texttt{raise exc}}
\begin{block}{Python / sample output}
\begin{Verbatim}[fontsize=\footnotesize]
try:
    parse_record(line)
except ValueError:
    log_bad_record(line)
    raise  # re-raise the same active exception
\end{Verbatim}
\end{block}
\small\begin{itemize}\item Bare \texttt{raise} inside an \texttt{except} re-raises the active exception.\item Use it when you record context but cannot recover.\item \texttt{raise exc} also raises it, but can alter traceback presentation.\end{itemize}
\end{frame}
\begin{frame}[fragile]{Add context without hiding the original cause}
\begin{block}{Python / sample output}
\begin{Verbatim}[fontsize=\footnotesize]
try:
    score = int(raw_score)
except ValueError as exc:
    raise ValueError("score must be an integer") from exc
\end{Verbatim}
\end{block}
\small\begin{itemize}\item \texttt{raise ... from exc} chains exceptions.\item The message gets more specific to your application.\item The original error remains available as the cause.\end{itemize}
\end{frame}
\begin{frame}[fragile]{Exception chaining, visualized}
\begin{center}\lectureimage{exception-chaining.png}\end{center}
\small\begin{itemize}\item Low-level failure: invalid integer.\item Higher-level failure: invalid score record.\item Both explanations can appear in the traceback.\end{itemize}
\end{frame}
\sectionframe{Part 3: Custom Exceptions}
\begin{frame}[fragile]{Why invent an exception type?}
\begin{itemize}\item \texttt{ValueError} tells us a value is unacceptable.\item But applications have specific problems: a malformed student-score record, for example.\item A custom exception lets callers catch \textbf{that kind of failure} separately.\item Start with a simple subclass of an existing built-in exception.\end{itemize}
\end{frame}
\begin{frame}[fragile]{Defining a custom exception is simple}
\begin{block}{Python / sample output}
\begin{Verbatim}[fontsize=\footnotesize]
class MalformedScoreError(ValueError):
    """A student score record is invalid."""
    pass

raise MalformedScoreError("expected name,score")
\end{Verbatim}
\end{block}
\small\begin{itemize}\item Subclassing \texttt{ValueError} means general \texttt{ValueError} handlers still work.\item The base class already stores the message; no custom \texttt{\_\_init\_\_} needed.\end{itemize}
\end{frame}
\begin{frame}[fragile]{Exception types form a family tree}
\begin{center}\lectureimage{exception-family.png}\end{center}
\small\begin{itemize}\item Built-in exceptions inherit from common parent classes.\item \texttt{MalformedScoreError} is also a \texttt{ValueError}.\item Use specific handler first if both are present.\end{itemize}
\end{frame}
\begin{frame}[fragile]{A parser with an explicit contract}
\begin{block}{Python / sample output}
\begin{Verbatim}[fontsize=\scriptsize]
def parse_score(line: str) -> tuple[str, int]:
    parts = line.strip().split(",")
    if len(parts) != 2:
        raise MalformedScoreError("expected name,score")
    name, raw = parts
    try:
        score = int(raw)
    except ValueError as exc:
        raise MalformedScoreError("bad integer") from exc
    if not 0 <= score <= 100:
        raise MalformedScoreError("score out of range")
    return name, score
\end{Verbatim}
\end{block}
\small\begin{itemize}\item Valid input returns a tuple. Invalid input raises.\item No silent fallback or magic number is returned.\end{itemize}
\end{frame}
\begin{frame}[fragile]{The caller handles the useful categories}
\begin{block}{Python / sample output}
\begin{Verbatim}[fontsize=\footnotesize]
try:
    with open(filename, encoding="utf-8") as f:
        records = [parse_score(line) for line in f]
except OSError as exc:
    print("Cannot read file:", exc)
except MalformedScoreError as exc:
    print("Bad score record:", exc)
\end{Verbatim}
\end{block}
\small\begin{itemize}\item I/O failures and malformed content are \textbf{different categories}.\item Neither exception needs to be caught inside \texttt{parse\_score()}.\end{itemize}
\end{frame}
\sectionframe{Part 4: Testing Exceptions}
\begin{frame}[fragile]{Two forms of \texttt{assertRaises}}
\begin{center}\lectureimage{assert-raises-two-forms.png}\end{center}
\small\begin{itemize}\item Function-call form: pass the function object and arguments.\item Context-manager form: put the failing call in the \texttt{with} body.\item Both pass only when the specified exception is raised.\end{itemize}
\end{frame}
\begin{frame}[fragile]{Syntax A: function-call form}
\begin{block}{Python / sample output}
\begin{Verbatim}[fontsize=\footnotesize]
import unittest

class ScoreTests(unittest.TestCase):
    def test_bad_score(self):
        self.assertRaises(
            MalformedScoreError, parse_score, "Ada,nope"
        )
\end{Verbatim}
\end{block}
\small\begin{itemize}\item Do \textbf{not} call \texttt{parse\_score(...)} before passing it.\item The test framework calls the function for you.\end{itemize}
\end{frame}
\begin{frame}[fragile]{Syntax B: context-manager form}
\begin{block}{Python / sample output}
\begin{Verbatim}[fontsize=\footnotesize]
class ScoreTests(unittest.TestCase):
    def test_bad_range(self):
        with self.assertRaises(MalformedScoreError):
            parse_score("Ada,150")
\end{Verbatim}
\end{block}
\small\begin{itemize}\item Best when the operation needs several lines of setup.\item Only place the expected failing statement inside the context.\end{itemize}
\end{frame}
\begin{frame}[fragile]{Check the message too}
\begin{block}{Python / sample output}
\begin{Verbatim}[fontsize=\footnotesize]
with self.assertRaisesRegex(
    MalformedScoreError, "out of range"
):
    parse_score("Ada,150")

with self.assertRaises(MalformedScoreError) as ctx:
    parse_score("Ada,150")
self.assertIn("range", str(ctx.exception))
\end{Verbatim}
\end{block}
\small\begin{itemize}\item \texttt{assertRaisesRegex} checks exception type and message.\item \texttt{as ctx} exposes the captured exception object.\end{itemize}
\end{frame}
\begin{frame}[fragile]{Common mistake: accidentally call it first}
\begin{columns}[T,onlytextwidth]
\begin{column}{0.48\textwidth}\begin{block}{Left}
\begin{Verbatim}[fontsize=\footnotesize]
WRONG
self.assertRaises(ValueError, int("oops"))
\end{Verbatim}
\end{block}\end{column}
\begin{column}{0.48\textwidth}\begin{block}{Right}
\begin{Verbatim}[fontsize=\footnotesize]
RIGHT
self.assertRaises(ValueError, int, "oops")

with self.assertRaises(ValueError):
    int("oops")
\end{Verbatim}
\end{block}\end{column}
\end{columns}
\vspace{0.06in}\begin{keyidea}\small The exception must be raised inside the test assertion, not before it.\end{keyidea}
\end{frame}
\begin{frame}[fragile]{Test the success path AND failure paths}
\begin{itemize}\item Valid: \texttt{parse\_score("Ada,87")} returns \texttt{("Ada", 87)}.\item Wrong shape: \texttt{"Ada"} raises \texttt{MalformedScoreError}.\item Wrong type: \texttt{"Ada,abc"} raises \texttt{MalformedScoreError}.\item Out of range: \texttt{"Ada,150"} raises \texttt{MalformedScoreError}.\item Use \texttt{assertEqual} for normal returns and \texttt{assertRaises} for exceptions.\end{itemize}
\end{frame}
\begin{frame}[fragile]{Mini-capstone: robust score importer}
\begin{center}\lectureimage{score-validation-pipeline.png}\end{center}
\small\begin{itemize}\item A CSV-style file contains lines like \texttt{"Ada,87"}.\item A malformed line should produce a \textbf{specific} error.\item The entry-point prints a useful message and exits cleanly.\end{itemize}
\end{frame}
\begin{frame}[fragile]{Walkthrough: assign responsibilities}
\begin{columns}[T,onlytextwidth]
\begin{column}{0.31\textwidth}\begin{block}{Parser}\small \texttt{parse\_score(line)}\\- validates one record\\- raises \texttt{MalformedScoreError}\end{block}\end{column}
\begin{column}{0.31\textwidth}\begin{block}{Loader}\small \texttt{load\_scores(path)}\\- opens and reads records\\- lets I/O errors propagate\end{block}\end{column}
\begin{column}{0.31\textwidth}\begin{block}{Entry point}\small \texttt{main()}\\- catches errors\\- prints user-facing feedback\end{block}\end{column}
\end{columns}
\vspace{0.06in}\begin{keyidea}\small Catch errors at the layer that can make a meaningful decision.\end{keyidea}
\end{frame}
\begin{frame}[fragile]{Your turn: choose a response}
\begin{activity}\small For each failure, should the helper return a value, raise an exception, or catch one?\\1. User enters "42" for a score.\\2. Score record contains "Ada,purple".\\3. File does not exist.\\4. A helper logs an error but cannot repair it.\end{activity}
\small\begin{itemize}\item Pair up: justify each choice using the function contract.\end{itemize}
\end{frame}
\begin{frame}[fragile]{Takeaways / exit ticket}
\begin{itemize}\item \texttt{try} / \texttt{except} handles \textbf{expected} exceptions; prefer specific types.\item \texttt{as exc} gives the raised object and its message.\item Uncaught exceptions propagate through function calls.\item \texttt{raise}, custom exception classes, and \texttt{raise ... from} communicate failures.\item Test failure behavior using \textbf{both} \texttt{assertRaises} \textbf{syntaxes}.\end{itemize}
\vspace{0.06in}\begin{keyidea}\small Exit ticket: When would you let \texttt{ValueError} propagate instead of catching it?\end{keyidea}
\end{frame}
\end{document}
